\documentclass[11pt]{article}
\usepackage[a4paper,margin=2.5cm]{geometry}
\usepackage{amsmath,amssymb,amsthm,booktabs}
\newtheorem{proposition}{Proposition}\theoremstyle{remark}\newtheorem*{remark}{Remark}
\title{The $(\log 2,\log 3]$ step:\\ \large Connes--Consani's archimedean operator rebuilt, the semi-local $\{\infty,2\}$ Sonin space on a slice,\\ and why the compact-remainder mechanism does not cross the place $2$}
\author{D. Joubert (with the assistance of two language models)}
\date{3 September 2026 --- consolidates notebook \S80--85 of \texttt{riemann-weil-phenomenology}}
\begin{document}\maketitle
\begin{abstract}
Weil positivity is known, without RH, for test functions supported in $[2^{-1/2},2^{1/2}]$ (window length $L=\log 2$); the proof of Connes and Consani (Selecta 2021) writes the archimedean Weil functional as a positive trace on Sonin's space minus a remainder whose conditioned operator is $-2\varepsilon'(1^+)\,\mathrm{Id}$ plus a compact part with a single eigenvalue above $1$. Crossing the threshold of the prime $2$ --- positivity for $L\in(\log 2,\log 3]$ --- is the first open step of their semi-local program. We (i) rebuild their archimedean operator and reproduce every published digit; (ii) measure that its compact part acquires a second eigenvalue above $1$ at $L\approx1.01$, while the archimedean form itself stays positive on the conditioned hyperplane at $L=\log 3$; (iii) show that the prime-$2$ term cannot be added to the compact remainder; (iv) build the semi-local $\{\infty,2\}$ Sonin space on the $\mathrm{ord}_2=0$ slice, where the adelic Fourier transform has the closed form $\mathfrak Fg(\rho)=\tfrac12[\sum_{n\ge0}\hat g(2^n\rho)-\hat g(\rho/2)]$, and find that the compression is not trace class, that the remainder is log-divergent at $\rho=1$ and spiked at $\rho=2$, and that the conditioned remainder is \emph{essentially positive} --- the sign that made the archimedean argument work is reversed; (v) check the construction against the semi-local trace formula of Connes (1999): the archimedean distribution is reproduced to $1$--$4\%$ and the $2$-adic place appears exactly at $\lambda=2^{\pm1}$, its weights not yet converged at $\Lambda\le8$. Everything is measurement and calibration; the conclusion is that the $(\log2,\log3]$ step requires a mechanism of a different nature, and the construction says which term is in the way. Nothing here bears on RH.
\end{abstract}

\section{The archimedean mechanism and where it stops}
Connes--Consani prove, for $g$ supported in $[2^{-1/2},2^{1/2}]$ with $\hat g(\pm i/2)=\hat g(0)=0$, that $W_\infty(g*g^*)\ge\mathrm{Tr}(\vartheta(g)S\vartheta(g)^*)\ge0$. The engine is the identity (their Theorem 4.7) $\mathrm{Tr}(\vartheta(f)S)=W_\infty(f)+E(f)$, $S$ the projection on Sonin's space $\mathcal S(1,1)$ (even functions vanishing with their Fourier transform on $[-1,1]$), $E(f)=\int f(\rho^{-1})\epsilon(\rho)\,d^*\rho$ with $\epsilon$ explicit in the prolate spheroidal functions of bandwidth $2\pi$. After the vanishing conditions are imposed through $Q=-(\rho\partial_\rho)^2+\tfrac14$, the remainder becomes $E\circ Q(\xi*\xi^*)=\langle\xi|N_I\xi\rangle$, $N_I=-2\epsilon'(1^+)(\mathrm{Id}-K_I)$, $K_I$ compact of Hilbert--Schmidt class on $L^2(I)$, $I=[-\tfrac12\log2,\tfrac12\log2]$. Positivity follows if $K_I\le1$; they compute that $K_I$ has exactly one eigenvalue above $1$ ($1.0516$), neutralized by one linear condition, with constant $c\in(13,17)$. The kink of $\epsilon$ at $\rho=1$ ($\epsilon'(1^+)=22.9965$) is what produces the $-2\epsilon'\mathrm{Id}$; the finiteness $\delta(1)=\sum\lambda(n)^2=2.2375$ is what makes the remainder a regular kernel. The whole argument stops at $L=\log 2$, to one eigenvalue.

\section{Reconstruction and calibration}
\texttt{code/cc\_arch.py} rebuilds the objects: prolates $PS_{2n,0}(2\pi,\cdot)$, eigenvalues $\lambda(n)$ of the finite Fourier transform, $\xi_n$, the analytic continuation $\xi_n^{an}(x)=\lambda(n)^{-1}\int_{-1}^1\xi_n(t)\cos(2\pi tx)\,dt$, $Q\epsilon$ from their formula (99), and $K_I$ by their $q\to1$ Toeplitz discretization ($\omega=2\times10^{-3}$).
\begin{center}\small
\begin{tabular}{lll}\toprule
quantity & rebuilt & Connes--Consani\\\midrule
$\lambda(0),\lambda(1),\lambda(2),\lambda(3)$ & $0.999971,\,-0.979485,\,0.524086,\,-0.058977$ & $0.999971,\,-0.979485,\,0.524086,\,-0.0589766$\\
$t(0),t(1),t(2),t(3)$ & $11.9719,\,8.77574,\,2.20528,\,0.0434$ & $11.9719,\,8.77574,\,2.20528,\,0.0433983$\\
$\epsilon'(1^+)$ & $22.9965$ & $22.9965$\\
$\lambda_{\max}(K_I),\lambda_2,\lambda_3$ at $\log 2$ & $1.05176,\,0.68791,\,0.0297$ & $1.05177\,(\omega=10^{-3}),\,0.687925,\,0.0289$\\\bottomrule
\end{tabular}
\end{center}

\section{Beyond $\log 2$ in the archimedean frame}
The second eigenvalue of $K_I$ crosses $1$ at $L\approx1.01$ ($\mu\approx2.74$): one eigenvalue above $1$ on $(\log2,1.01)$, two on $(1.01,\log3]$ ($1.0899$, $1.0395$ at $\log3$). Yet the archimedean form itself is still positive in their frame: on the Galerkin space $V_{31}$ at $\mu=3$, the archimedean form with the pole removed has one negative direction on $V$ (along the pole) and \emph{none} on the hyperplane $c(f)=0$ of the vanishing condition. Their sufficient condition $K_I\le1$ thus already fails at $\log3$ while the conclusion holds --- the Sonin trace compensates both directions of the remainder.

The prime-$2$ term cannot be added to the remainder. On the conditioned functions $\xi$, $W_2\circ Q(\xi*\xi^*)=\tfrac{\log2}{\sqrt2}\bigl[\tfrac14F(\log2)-F''(\log2)\bigr]$, $F=\xi*\xi^*$, is a differential form, unbounded on $L^2(I)$; on the side $g=(\tfrac12-\partial)\xi$ it is a bounded translation form, but the archimedean identity term becomes $-2\epsilon'\|Y*g\|^2$, a weak norm that does not control $G(\log2)$ for oscillating $g$. In either presentation $E+W_2\le0$ is false, although $W_\infty-W_2\ge0$ holds numerically at $\log3$ ($\lambda_{\min}=5.9\times10^{-8}$ on $V_{31}$): the positivity comes from the positive trace absorbing $W_2$, not from a bound on the remainder. This is the operator form of a Galerkin measurement: at $L=\log3$ the archimedean form $Q_\infty=\mathrm{pole}+\mathrm{arch}$ has exactly one negative direction $v$ (eigenvalue $-0.074$), $-T_2$ is positive on it with margin ($Q(v)=+0.050$, using $5\%$ of $T_2$'s mass), and the razor $5.9\times10^{-8}$ lives in the positive complement.

\section{The semi-local $\{\infty,2\}$ Sonin space on a slice}
Following Connes (1999, \S VII), $X_S=(\mathbb R\times\mathbb Q_2)/\pm2^{\mathbb Z}$, $L^2(X_S)$ is the completion of $\mathcal S(\mathbb A_S)$ for $\|f\|^2=\int_{\mathcal D}|\sum_q f(qx)|^2|x|\,d^*x$, and $F=F_\infty\otimes F_2$ is unitary on it (his Lemma 1). Functions invariant under $\mathbb Z_2^*$ are determined by their restriction $g$ to the shell $\mathrm{ord}_2=0$, and on this slice
\begin{equation}\label{F}
\mathfrak Fg(\rho)=\tfrac12\Bigl[\sum_{n\ge0}\hat g(2^n\rho)-\hat g(\rho/2)\Bigr],\qquad \hat g(\xi)=2\int_0^\infty g(r)\cos(2\pi r\xi)\,dr .
\end{equation}
\begin{proposition}
$\mathfrak F$ is unitary and involutive on $L^2([0,\infty))$; hence self-adjoint.
\end{proposition}
\begin{proof}
Write $\mathfrak Fg=\sum_{k\ge-1}c_k\,\hat g(2^k\cdot)$ with $c_{-1}=-\tfrac12$, $c_k=\tfrac12$ ($k\ge0$). For $j\ne0$ the cross coefficient $\sum_k c_kc_{k+j}2^{-k}$ vanishes ($-\tfrac12+\tfrac12$), and the diagonal $\sum_kc_k^22^{-k}=\tfrac12+\tfrac12=1$: $\|\mathfrak Fg\|=\|\hat g\|=\|g\|$. Involutivity is $F^2=\mathrm{id}$ on even functions of $\mathbb A_S$, which descends. Numerically (\texttt{code/semilocal.py}, tests): $\|\mathfrak Fg\|^2/\|g\|^2=1.015$, $\mathfrak F^2g=0.94\,g$ on the support of a Gaussian bump, the deviations being quadrature.
\end{proof}
The cutoff $P_1=\mathbf 1_{[0,1]}$ and the scaling $\vartheta(\lambda)g(r)=\lambda^{-1/2}g(r/\lambda)$ are unchanged; only $\hat P_1=\mathfrak FP_1\mathfrak F$ and Sonin's space change. Three facts, all computed with exact cell averaging (a first pointwise assembly aliased the lacunary terms and produced a spurious $43\%$ asymmetry; it is retracted in the notebook):
\begin{itemize}
\item \emph{The compression is not trace class.} The eigenvalues of $P_1\mathfrak FP_1$ decay slowly ($0.998,-0.990,0.613,-0.554,0.429,0.407,\ldots$) and $\sum\lambda_n^2$ grows with resolution ($4.24$ at $N=200$, $4.75$ at $N=400$; archimedean: $2.2375$, stable). The kernel $\sum_n\cos(2\pi2^n\rho r)$ is not square integrable on $[0,1]^2$.
\item \emph{The remainder is log-divergent at $\rho=1$ and spiked at $\rho=2$.} $\delta_S(\rho)=\mathrm{Tr}(\vartheta(\rho^{-1})\hat P_1P_1)=\sum_n\lambda_n\langle\xi_n|\vartheta(\rho^{-1})\eta_n\rangle$ (CC's Lemma 4.1, which uses only the algebra of two projections). Archimedean control: the closed form $2\sqrt\rho\,[\mathrm{Si}(2\pi(1+\rho))/(2\pi(1+\rho))+\mathrm{Si}(2\pi(\rho-1))/(2\pi(\rho-1))]$ is reproduced to $3$--$4$ digits. Semi-local: $\delta_S\approx-0.65\ln(\rho-1)+C$ ($\delta_S(1)=\sum\lambda_n^2=\infty$) --- a logarithm where CC had a kink --- and a spike at $\rho=2$ growing with resolution: the $2$-adic Weil term, appearing where it must.
\item \emph{The conditioned remainder is essentially positive.} On twenty smooth test functions of $I$ (sine basis; eigenvalues of $D\circ Q$ relative to the Gram matrix; \texttt{code/dq\_sign.py}): archimedean at $\log2$, $2$ positive eigenvalues out of $20$ ($4.77$, $1.23$) and the rest massed at $-2.01$ --- the $-2\,\mathrm{Id}$ of CC's Theorem 3.6 in plain sight, and their Remark 3.9 recovered; at $\log3$, $3$ positive. Semi-local at $\log2$: \textbf{$20$ positive out of $20$} (up to $121$); at $\log3$, $13$ out of $20$ (up to $196$), the positive part growing with the basis --- the $H^{1/2}$ signature of the logarithm ($-\ln|x|$ and the finite part of $1/x^2$ both have positive transforms).
\end{itemize}

\section{The verdict}
$W_\infty-W_2=L_S-D_S$ with $L_S$ a positive trace. Extracting positivity requires $D_S\circ Q$ essentially negative, as in the archimedean case. It is essentially positive: the positive trace absorbs the place $2$ but with a divergent positive surplus that erases the sign of $W$. The compact-remainder mechanism of Connes--Consani, transported verbatim to $\{\infty,2\}$ at $\Lambda=1$, does not give the $(\log2,\log3]$ step (28th preregistered execution of the repository). The construction also says why: the surplus is the term $2h(1)\log'\Lambda$ of the semi-local trace formula of Connes (1999), and that theorem states that the finite part of the compressed trace, once this term is subtracted, \emph{is} the Weil functional --- no positivity comes for free. The archimedean miracle (a finite $\delta(1)$ and a kink) belongs to one place.

\section{Validation against the semi-local trace formula}
$\tau_\Lambda(\lambda)=\mathrm{Tr}(\hat P_\Lambda P_\Lambda\vartheta(\lambda))$ on the slice (\texttt{code/trace\_dist.py}, $\Lambda=4,8$). Theorem 4 of Connes (1999) predicts $\tau_\Lambda\to2\log'\Lambda\,\delta_1+\tau_\infty+\tau_2$ with $\tau_\infty=\lambda^{1/2}/2\,(1/(1+\lambda)+1/|1-\lambda|)$ (CC (39)) and $\tau_2$ point masses at $\lambda=2^{\pm m}$. Archimedean: $\tau_\Lambda$ matches (39) to $1$--$4\%$ away from $\lambda=1$ ($0.7495/0.7529$ at $0.4$, $2.4986/2.4845$ at $0.8$, $0.9338/0.9428$ at $2.0$, $0.6521/0.6495$ at $3.0$). Semi-local minus archimedean: two sharp peaks exactly at $\lambda=2$ and $\lambda=\tfrac12$ at both $\Lambda$ --- the $2$-adic place where the theorem puts it. Their integrated weights ($+0.13/-0.08$ at $\Lambda=4$, $-0.23/-0.35$ at $\Lambda=8$, against $\log2/\sqrt2=0.49$) oscillate and change sign: not converged at $\Lambda\le8$. The check is conclusive on the location of the places and open on the $2$-adic weight; closing it needs $\Lambda\ge16$ (matrices $512\times1700$ with forty lacunary Si terms) or an analytic treatment of the $2$-adic shells.

\section*{Status and reproducibility}
Calibration is exact to the published digits; the semi-local objects are computed with exact cell averaging and covered by fourteen recomputing tests (\texttt{tests/test\_semilocal\_*.py}, suite of $128$). The sign reversal of the conditioned remainder is a measurement on twenty test functions at two window lengths, not a theorem; the identification of the surplus with the $\log\Lambda$ term of 1999 is an inference, supported by the log-divergence of $\delta_S$ at $1$. What the note establishes is negative and precise: the first semi-local step of the Connes--Consani program cannot be taken by transporting their $\Lambda=1$ Sonin mechanism; it needs a term subtracted before the compact part can be examined. Nothing here bears on RH.
\end{document}
